\section{Introduction}\label{sec:intro}
Infrared and visible sensors capture different aspects of a scene. Visible images describe reflected-light appearance and texture, but their quality depends on illumination. Infrared images reveal thermal patterns that remain useful in low-light environments, although they often contain less appearance detail. Infrared and visible image fusion (IVIF) combines these complementary observations into a single image, with applications in nighttime perception and scene monitoring~\cite{ma2019survey,tang2022seafusion}. An effective method must preserve thermal targets and visible structures while keeping computation manageable.

Deep learning has advanced this task through convolutional feature extraction, auto-encoding, adversarial learning, and attention-based modeling~\cite{li2019densefuse,xu2022u2fusion,ma2019fusiongan,ma2022swinfusion}. Feature-decomposition methods provide a useful way to organize the information to be fused. DIDFuse~\cite{zhao2020didfuse} separates background and detail representations, and CDDFuse~\cite{zhao2023cddfuse} uses a correlation-driven dual-branch design to encourage shared base features and complementary details. Its reconstruction-then-fusion training provides a strong foundation, but leaves room to improve how computation is allocated and how fusion quality is supervised.

We focus on three connected issues. The first is the cost of full-resolution processing. CDDFuse uses Restormer-style blocks for shallow encoding and reconstruction. Their transposed attention efficiently models channel relationships without constructing a quadratic spatial attention matrix~\cite{zamir2022restormer}; nevertheless, repeated projections and feed-forward transformations still contribute substantial computation. Since local boundaries and textures are central to these stages, lightweight gated convolutions offer a way to reduce cost while retaining richer context modeling in the base branch.

The second issue is combining spatial context with channel relationships. Transposed attention derives channel interactions from features aggregated over the image. Selective state-space models offer a complementary route by propagating information through spatial sequences~\cite{gu2023mamba,liu2024vmamba}. Recent fusion methods have explored Mamba for feature enhancement and cross-modal interaction~\cite{xie2024fusionmamba,li2024mambadfuse}. For dense fusion, however, context aggregation should coexist with a direct path for full-resolution features. This motivates a pooled scanning path anchored to a full-resolution residual stream, followed by channel attention. Difference-guided fusion (DGF) and magnitude-guided fusion (MGF) then use cross-modal differences and feature magnitudes to adjust base and detail aggregation, respectively.

The third issue is the balance between intensity, edges, structure, and contrast. Maximum-intensity and gradient targets give useful supervision, but do not explicitly regularize normalized local correlation or global and local intensity variation. The lack of a unique fused ground truth makes this balance consequential: different objectives emphasize different image properties~\cite{zhao2025freefusion}. We retain the original intensity and gradient targets and add complementary structure and contrast terms. An intra-modal decorrelation penalty also discourages redundant patterns between the learned base and detail representations.

We therefore propose SCMFuse, an efficient space--channel Mamba fusion framework. It combines multi-scale local processing, hybrid spatial--channel modeling, adaptive feature aggregation, and targeted supervision within the CDDFuse decomposition paradigm. The main contributions are:
\begin{itemize}
 \item \textbf{Efficient local processing.} A shallow gated convolution block (SGC Block) and a hybrid decoder reduce the computation used for full-resolution feature extraction and reconstruction.
 \item \textbf{Spatial--channel modeling with adaptive fusion.} A space--channel MambaVision block (SCM Block) integrates four-directional spatial scanning, full-resolution residual features, and channel attention. Difference- and magnitude-guided fusion adapt the aggregation of base and detail features.
 \item \textbf{Complementary training constraints.} Intra-modal decorrelation and contrast--structure regularization complement the baseline decomposition and fusion objectives. Experiments on three datasets demonstrate improvements in structural correlation and cross-dataset edge fidelity alongside substantial model compression and acceleration.
\end{itemize}
The full model uses $0.473$ M parameters and $87.837$ G FLOPs, reducing the corresponding CDDFuse profile by $60.2\%$ and $62.7\%$. Its $2.76\times$ mean inference speedup and improved TNO/RoadScene fidelity illustrate the benefit of allocating computation according to the roles of local extraction, global context, and feature aggregation.
